\appendix

\section{Experimental Apparatus}
\label{app:apparatus}

Figure~\ref{fig:pipeline} shows the recursive procedure. One \emph{chain} is one
independently seeded traversal of the loop for the full horizon $G$. Chains, not
generations, are the experimental units: generations within a chain are repeated
observations of the same trajectory and are not independent samples.

\begin{figure}[h]
\centering
% Recursive train/generate/allocate/evaluate pipeline.
% Result-independent: this figure describes the experimental apparatus, not any
% outcome, so it is safe to include before results exist.
% Requires \usepackage{tikz} and \usetikzlibrary{arrows.meta,positioning} in main.tex.
\begin{tikzpicture}[
  node distance=9mm and 12mm,
  every node/.style={font=\small},
  stage/.style={
    draw, rounded corners=2pt, align=center,
    minimum height=8mm, minimum width=22mm, inner sep=4pt
  },
  frozen/.style={stage, fill=black!5},
  artifact/.style={stage, densely dashed},
  flow/.style={-{Stealth[length=2mm]}, thick},
  feedback/.style={-{Stealth[length=2mm]}, thick, densely dotted},
]

  \node[frozen] (manifests) {Frozen data\\manifests};
  \node[stage, right=of manifests] (train) {Train\\generation $g$};
  \node[stage, right=of train] (generate) {Generate\\synthetic corpus};
  \node[stage, right=of generate] (allocate) {Allocate human\\rescue $b_{g,m}$};

  \node[frozen, below=of manifests] (budget) {Lifetime budget $B$\\$\sum_{g,m} b_{g,m} = B$};
  \node[stage, below=of generate] (evaluate) {Held-out\\evaluation};
  \node[artifact, below=of evaluate] (chain) {Immutable chain\\result + manifest};
  \node[artifact, right=of chain] (analysis) {Chain-level\\statistics};

  \draw[flow] (manifests) -- (train);
  \draw[flow] (train) -- (generate);
  \draw[flow] (generate) -- (allocate);

  % The recursion: this generation's allocation feeds the next generation's training.
  \draw[feedback] (allocate.north) -- ++(0,6mm) -| node[pos=0.25, above] {generation $g{+}1$} (train.north);

  \draw[flow] (budget) -- (allocate.south -| budget.east) ;
  \draw[flow] (budget.east) -- (evaluate.west);
  \draw[flow] (generate.south) -- (evaluate.north);
  \draw[flow] (evaluate) -- (chain);
  \draw[flow] (chain) -- (analysis);

\end{tikzpicture}

\caption{Recursive train--generate--allocate--evaluate pipeline. Shaded boxes are
frozen before any run. Dashed boxes are immutable artifacts. The dotted edge is
the recursion: generation $g$'s allocation enters generation $g{+}1$'s training
mixture. The lifetime budget constrains every allocation jointly, so spending
early reduces what remains for later generations.}
\label{fig:pipeline}
\end{figure}

\section{Data Appendix}
\label{app:data}

\textbf{RESULT\_PENDING.} The primary domain is not yet frozen
(\texttt{DECISIONS.md} U-002). This appendix records what will be reported once it
is, and is structured so that filling it requires no reorganisation.

\begin{itemize}
  \item Dataset identifier, revision hash, and licence, with the licence audit
    that selected it over the audited alternatives.
  \item Preprocessing: tokenizer identity and revision, normalisation, and the
    exact filter chain, each with a content hash.
  \item Partition statistics for all five partitions --- base training, rescue
    candidates, prompts, validation, and final test --- reporting example counts,
    token counts, and per-partition manifest hashes.
  \item The frozen mode definition, the number of modes $M$, and the criterion
    separating tail modes from common modes.
  \item Overlap report: exact and near-duplicate checks across every partition
    pair, with the near-duplicate threshold declared before it was applied.
\end{itemize}

Final test data influences no prompt, selection decision, threshold, stopping
rule, or hyperparameter. This constraint is enforced by the disjointness tests
rather than asserted.

\section{Evaluation Appendix}
\label{app:evaluation}

\textbf{RESULT\_PENDING.} The primary tail-retention definition is not yet frozen
(U-004). Two candidates are implemented and audited for determinism,
independence from the policy's selection signal, sensitivity to tail degradation,
discriminant validity against non-tail modes, and split-half reliability.

Held-out human negative log-likelihood is computed from target-token log
probabilities on the frozen test partition. Regret is aggregated across the
horizon as the area under the generation-wise curve.

Both tail candidates are computed without reference to the policy's
under-coverage score, so no policy can be rewarded for agreeing with its own
selection signal. This is a structural property of the interfaces, not a
convention.

\section{Statistics Appendix}
\label{app:statistics}

\textbf{RESULT\_PENDING.} Uncertainty is computed across independent chains.
Multiple generations within one chain are repeated observations and do not
increase the effective sample size.

The primary contrast is the joint policy against the strongest eligible non-joint
baseline, paired by chain seed. Paired differences are summarised by their mean
with a bootstrap interval resampled over complete chains. The smallest
scientifically meaningful effect is frozen before outcomes are opened (U-006),
and the interpretation that follows is selected mechanically from precommitted
templates covering favourable, null, harmful, and uncertain outcomes.

\section{Reproducibility Appendix}
\label{app:reproducibility}

Every value reported in this paper is generated by a script and traced to an
immutable artifact. No number is entered by hand.

\begin{itemize}
  \item \textbf{Positive control.} Upstream repository and exact commit are
    pinned before execution, together with the resolved framework versions and
    tokenizer revision actually used. Deviations from the upstream configuration
    are declared before running, never after observing a result.
  \item \textbf{Run artifacts.} Each chain emits a run manifest and a chain
    result validated against versioned JSON schemas. A three-state validator
    classifies each run as valid, valid with limitation, or invalid; blocking
    failures include budget-equality violations, identity mismatches, and
    impossible token accounting.
  \item \textbf{Budget matching.} Every policy in a comparison consumes identical
    lifetime human-origin optimizer tokens and identical total optimizer tokens.
    Equality is checked before launch from configuration and again after the run
    from emitted artifacts.
  \item \textbf{Token accounting.} Counts derive from tokenized batches actually
    consumed by the optimizer. Repeated presentation of an example costs its
    tokens once per presentation; gradient accumulation counts every micro-batch;
    resume neither double-counts nor drops the batch spanning a checkpoint.
  \item \textbf{Determinism.} Seeds propagate through sampling, generation,
    initialisation, dropout, and evaluation. Checkpoints carry content digests so
    silent corruption cannot survive a resume.
  \item \textbf{Generated outputs.} Tables and figures are emitted by scripts and
    recorded with SHA-256 digests in an artifact manifest, so a reader can verify
    that a figure is the one the pipeline produced.
\end{itemize}

\section{Limitations of Scope}
\label{app:scope}

\textbf{RESULT\_PENDING.} The planned study covers one licensed domain, one
124M--160M screening model, and one horizon. It cannot establish behaviour at
larger scales, in other domains, or for other model families. Any claim beyond
the tested configuration is future work and is marked as such.
